\section{Conjunto de teste com anomalia}

Esta seção aplica o modelo identificado nos dados normais ao conjunto
\texttt{dados\_teste\_anomalia.csv}, calcula os resíduos e os analisa para detectar
anomalias.

\subsection{Aplicação do modelo}

O modelo FIR obtido no conjunto de entrada branca foi aplicado à entrada do conjunto
de teste, gerando a saída prevista. A Figura~\ref{fig:C_sinais} mostra a entrada, a
saída verdadeira e a saída prevista. Na primeira parte do registro, as duas saídas se
acompanham. A partir de certo ponto, a saída verdadeira passa a apresentar picos que
a previsão não reproduz.

\begin{figure}[H]
\centering
\includegraphics[width=0.95\textwidth]{C_sinais.png}
\caption{Entrada, saída verdadeira e saída prevista no conjunto de teste.}
\label{fig:C_sinais}
\end{figure}

\subsection{Cálculo dos resíduos}

Os resíduos foram calculados como $e[n] = y[n] - \hat{y}[n]$, descartando as
primeiras amostras correspondentes ao transitório do filtro. Como referência de
normalidade, usou-se o desvio padrão dos resíduos da validação com entrada branca,
igual a \valCbaselineStd, faixa em que o modelo é válido e não há anomalia. A
Figura~\ref{fig:C_residuos} mostra os resíduos ao longo do tempo.

\begin{figure}[H]
\centering
\includegraphics[width=0.9\textwidth]{C_residuos.png}
\caption{Resíduos $e[n]$ no conjunto de teste. As linhas tracejadas marcam três vezes
o valor de referência.}
\label{fig:C_residuos}
\end{figure}

\subsection{Análise dos resíduos}

\subsubsection{Estatísticas globais}

A média dos resíduos permanece próxima de zero nos dois regimes
(\valCmediaENormal{} no trecho normal e \valCmediaEAnom{} no trecho anômalo), o
que indica que a anomalia não introduz um desvio sistemático de nível: ela se
manifesta na variância e na dinâmica. Já a variância global dos resíduos do
conjunto de teste é cerca de \valCvarRatio{} vezes maior do que a de
referência, o que já aponta uma diferença de comportamento em relação ao caso
normal.

\subsubsection{Energia local}

A energia local, definida como a média móvel do quadrado dos resíduos, está na
Figura~\ref{fig:C_energia_local}. Ela permanece próxima do valor de referência na
primeira parte do registro e em seguida sobe, chegando a um pico de cerca de
\valCenergiaRatio{} vezes a referência. Esse salto marca uma mudança de regime em
torno de $n \approx \valCnChange$.

\begin{figure}[H]
\centering
\includegraphics[width=0.9\textwidth]{C_energia_local.png}
\caption{Energia local dos resíduos no conjunto de teste.}
\label{fig:C_energia_local}
\end{figure}

\subsubsection{Autocorrelação}

A Figura~\ref{fig:C_autocorr_residuos} mostra a autocorrelação dos resíduos para o
registro completo e para os dois regimes separados pela mudança detectada. No regime
normal, a fração de atrasos fora da banda é de \valCautoNormal, ou seja, os resíduos
são próximos de um ruído branco. No regime anômalo, essa fração sobe para
\valCautoAnom, o que indica que o modelo deixou de descrever a dinâmica naquele
trecho.

\begin{figure}[H]
\centering
\includegraphics[width=0.98\textwidth]{C_autocorr_residuos.png}
\caption{Autocorrelação dos resíduos: registro completo, regime normal e regime
anômalo.}
\label{fig:C_autocorr_residuos}
\end{figure}

\subsection{Detecção de anomalias}

Foram combinados três indicadores simples sobre os resíduos: a energia local, para
mudança de regime; a variância por blocos, para aumento do erro; e um limiar de
amplitude, para impulsos. A Figura~\ref{fig:C_deteccao} reúne a variância por blocos e
a marcação dos impulsos. A Tabela~\ref{tab:C_deteccao} resume os números.

\begin{figure}[H]
\centering
\includegraphics[width=0.95\textwidth]{C_deteccao.png}
\caption{Variância dos resíduos por bloco e impulsos detectados no conjunto de teste.}
\label{fig:C_deteccao}
\end{figure}

\begin{table}[H]
\centering
\caption{Resumo da detecção nos dois regimes.}
\label{tab:C_deteccao}
\begin{tabular}{lcc}
\toprule
Indicador & Regime normal & Regime anômalo \\
\midrule
Desvio padrão dos resíduos & \valCstdNormal & \valCstdAnom \\
Razão em relação à referência & --- & \valCstdAnomRatio \\
Maior valor absoluto & --- & \valCmaxAnom \\
\bottomrule
\end{tabular}
\end{table}

A mudança de regime foi localizada em torno de $n \approx \valCnChange$. No regime
anômalo, o desvio padrão dos resíduos é cerca de \valCstdAnomRatio{} vezes o valor de
referência. Foram detectados \valCnumImpulsos{} impulsos com amplitude acima de oito
vezes a referência.

\subsection{A anomalia no sinal de saída ou nos resíduos}

Resta comparar onde a anomalia aparece com mais clareza. A
Figura~\ref{fig:C_comparacao} mostra, em cima, a energia local do sinal de saída e,
embaixo, a energia local dos resíduos. A Tabela~\ref{tab:C_comparacao} reúne algumas
medidas para os dois sinais.

\begin{figure}[H]
\centering
\includegraphics[width=0.95\textwidth]{C_comparacao_y_residuo.png}
\caption{Energia local do sinal de saída e dos resíduos no conjunto de teste.}
\label{fig:C_comparacao}
\end{figure}

\begin{table}[H]
\centering
\caption{Comparação entre o sinal de saída e os resíduos. Os picos são medidos
no regime anômalo em unidades do desvio padrão do regime normal de cada sinal.}
\label{tab:C_comparacao}
\begin{tabular}{lcc}
\toprule
Medida & Saída $y[n]$ & Resíduos $e[n]$ \\
\midrule
Razão de variância (anômalo/normal) & \valCvarRatioY & \valCvarRatioE \\
Pico em relação ao desvio normal & \valCpicoStdY & \valCpicoStdE \\
Atrasos de autocorrelação fora da banda (normal) & \valCautoNormalY & \valCautoNormal \\
Atrasos de autocorrelação fora da banda (anômalo) & \valCautoAnomY & \valCautoAnom \\
\bottomrule
\end{tabular}
\end{table}

A anomalia aparece nos dois sinais, mas fica mais clara nos resíduos. O sinal de saída
sempre tem variabilidade, porque é a resposta de um sistema excitado por ruído, e por
isso a mudança eleva a sua variância apenas cerca de \valCvarRatioY{} vezes. Nos
resíduos, que removem a parte previsível da saída, o contraste é bem maior: a
variância do regime anômalo é cerca de \valCvarRatioE{} vezes a do regime normal e os
impulsos se destacam de um patamar próximo de zero. Aplicar o mesmo limiar de impulso
ao sinal de saída marcaria \valCimpulsosY{} amostras, contra \valCnumImpulsos{} nos
resíduos, porque a saída atinge amplitudes altas mesmo sem anomalia. A autocorrelação
conta a mesma história: em $y[n]$ ela já está fora da banda no regime normal
(fração \valCautoNormalY), por se tratar de um sinal filtrado, e por isso não
discrimina nada; nos resíduos ela passa de \valCautoNormal{} para \valCautoAnom{} e
revela que o modelo deixou de explicar a dinâmica. Por esses motivos, a análise dos
resíduos é o caminho mais adequado para a detecção.
